\documentclass[11pt]{amsart}
\usepackage[T1]{fontenc}
\usepackage{amsmath,amssymb,amsthm}
\usepackage{booktabs}
\usepackage[hidelinks]{hyperref}
\hypersetup{pdfauthor={Daniele Cappello and the UNICO/NOUS pipeline},
  pdftitle={The golden moment and certified constraints on Agrawal's conjecture at r=5}}

\newtheorem{theorem}{Theorem}[section]
\newtheorem{lemma}[theorem]{Lemma}
\newtheorem{corollary}[theorem]{Corollary}
\newtheorem{conjecture}[theorem]{Conjecture}
\theoremstyle{remark}
\newtheorem{remark}[theorem]{Remark}

\newcommand{\eps}{\varepsilon}
\newcommand{\Z}{\mathbf{Z}}
\newcommand{\F}{\mathbf{F}}
\newcommand{\lean}[1]{\texttt{#1}}
\emergencystretch=1.5em
\raggedbottom

\begin{document}

\title[The golden moment at \(r=5\)]{The golden moment and certified
constraints on Agrawal's conjecture at \(r=5\)}

\author{Daniele Cappello and the UNICO/NOUS pipeline}
\thanks{This is a human-directed, LLM-assisted project. The mathematical
exploration and formalization were produced by an orchestrated multi-model
pipeline (Claude and GPT-family models), directed by Daniele Cappello and
subjected to independent model audits. Every formal claim is checked by the
Lean kernel; model agreement is not treated as evidence. Repository with the
Lean sources and all computational certificates:
\url{https://github.com/Solarys431/agrawal-r5}.}

\begin{abstract}
We prove, and formally verify in Lean, that the quadratic moment unit
arising in Agrawal's congruence at \(r=5\) has quintic index
\[
M_2=3\,\operatorname{ind}_5\!\left(\frac{1+\sqrt5}{2}\right).
\]
Thus the quadratic moment is exactly the classical quintic character of
the golden unit, up to the scalar \(3\). This note also collects the
mechanically verified constraints obtained around that bridge. Every
statement called a theorem below is proved in Lean~4 over Mathlib, with no
\lean{sorry}, no \lean{admit} and no added axioms; the corresponding
declaration name is given with each statement. Every computational claim
carries a replayable certificate with its detection criterion embedded and
SHA-256 manifests. We pose, without any claim of proof, the two open
problems that remain central. Nothing in this note asserts Agrawal's
conjecture.
\end{abstract}

\maketitle
\enlargethispage{4pt}

\section{Setting}

Agrawal's conjecture (arising from the AKS primality test
\cite{AKS}; see also \cite{LPremarks,LPgauss,Vana}) is stated for coprime
positive integers \(n,r\). We use its prime-\(r\) specialization: if \(r\)
is prime, \(r\nmid n\), and
\((X-1)^n\equiv X^n-1 \pmod{n,\,X^r-1}\),
then \(n\) is prime or \(n^2\equiv1\pmod r\). Our study concerns
\(r=5\). This note isolates its mechanically certified core; the
statements are elementary, and their interest lies in the role they
play in the wider study and in the fact that a proof assistant has
checked every one of them.

Local notation and all parity hypotheses are introduced at the point of
use. In the first sections, let
\[
R_p \;=\; (\Z/p)[X]\,/\,(X^2-X-1),
\qquad
\eps \;=\; \text{the class of } X,
\]
so that \(\eps^2=\eps+1\), and set \(\sqrt5:=2\eps-1\), which
satisfies \((\sqrt5)^2=5\) in \(R_p\)
(\lean{GoldenRing}, \lean{eps\_sq}, \lean{sqrt5\_sq}).
Let \(F_n\), \(L_n\) denote the Fibonacci and Lucas numbers, with
\(L_n=F_{n+1}+F_{n-1}\) (\lean{lucas}). Define
\[
A_n=\begin{cases}L_n,& n \text{ even},\\ F_n,& n\text{ odd},\end{cases}
\qquad
H_n=\gcd\bigl(A_n,\,5^{\,n-1}+1\bigr),
\]
and, dually, \(B_n=F_n\) (\(n\) even), \(L_n\) (\(n\) odd), with
\(J_n=\gcd\bigl(B_n,\,5^{\,n-1}-1\bigr)\).

\section{The golden Frobenius and the inertia of \texorpdfstring{\(J_n\)}{Jn}}

\begin{theorem}[\lean{golden\_frobenius}]\label{thm:frob}
Assume \(5\) is not a square in \(\Z/p\) and \(p\neq2\). Then
\(\eps^p=1-\eps\) in \(R_p\).
\end{theorem}

\begin{proof}[Certified proof sketch]
By Euler's criterion (\lean{euler\_nonsquare}) we have
\(5^{(p-1)/2}=-1\) in \(\Z/p\), hence
\[
(\sqrt5)^p=(5)^{(p-1)/2}\sqrt5=-\sqrt5
\qquad(\lean{sqrt5\_pow\_p}).
\]
Expanding \((2\eps-1)^p\) by the Frobenius endomorphism and
cancelling the unit \(2\) gives the claim.
\end{proof}

\begin{corollary}[\lean{golden\_pow\_p\_succ}]\label{cor:halfperiod}
Under the same hypotheses, \(\eps^{\,p+1}=-1\) in \(R_p\).
\end{corollary}

\begin{theorem}[\lean{inertia\_J}]\label{thm:inertia}
Assume \(5\) is not a square in \(\Z/p\) and \(p\neq2\). Then for no
\(n\ge1\) do \(\eps^{2n}=1\) in \(R_p\) and \(5^{\,n-1}=1\) in
\(\Z/p\) hold simultaneously.
\end{theorem}

\begin{proof}[Certified proof sketch]
If \(n\) is odd, raising \(\eps^{2n}=1\) to the power \((p+1)/2\) and
comparing with \((\eps^{p+1})^n=(-1)^n=-1\) gives \(1=-1\), absurd for
\(p\neq2\). If \(n\) is even, \(5=5^{\,n}=(5^{\,n/2})^2\) exhibits
\(5\) as a square, against the hypothesis.
\end{proof}

Via the identity \(\eps^{\,n+1}=F_{n+1}\eps+F_n\) in \(R_p\)
(\lean{golden\_pow\_fib} and its predecessor form
\lean{golden\_pow\_pred}) one obtains the divisibility forms, in
which the hypothesis is purely arithmetic by quadratic reciprocity
(\lean{not\_isSquare\_five}):

\begin{theorem}[\lean{inertia\_J\_fib\_mod5},
\lean{inertia\_J\_lucas\_mod5}]\label{thm:inertiadvd}
Let \(p\equiv\pm2\pmod5\), \(p\neq2\). For even \(n\ge2\), \(p\) does
not divide \(\gcd(F_n,\,5^{\,n-1}-1)\); for odd \(n\ge1\), \(p\) does
not divide \(\gcd(L_n,\,5^{\,n-1}-1)\). Equivalently, no \emph{odd}
prime inert in \(\mathbf{Q}(\sqrt5)\) divides any \(J_n\).
\end{theorem}

\section{The canonical support witness}

\begin{theorem}[\lean{support\_witness}]\label{thm:support}
Let \(p\equiv2\pmod5\), \(p\neq2\), and put \(n_0=(p+1)/2\). Then
\(n_0\equiv4\pmod5\), \(p\mid 5^{\,n_0-1}+1\), and \(p\) divides
\(L_{n_0}\) if \(n_0\) is even, \(F_{n_0}\) if \(n_0\) is odd. In
particular \(p\mid H_{n_0}\), with an index in the class
\(4\bmod5\).
\end{theorem}

\section{The moment obstruction}\label{sec:moment-obstruction}

Let \(F\) be a finite field, let \(e:F^\times\to F\), and define
\[
M_j(e)=\sum_{a\in F^\times}a^j e(a).
\]
The discrete logarithms occurring in the local Agrawal problem satisfy a
covariance relation of the following form.

\begin{theorem}[\lean{moment\_covariance}]\label{thm:moment-covariance}
If \(t\in F^\times\) and
\[
e(ta)=t\,e(a)\qquad(a\in F^\times),
\]
then
\[
tM_j(e)=t^{-j}M_j(e).
\]
\end{theorem}

\begin{proof}[Certified proof sketch]
Multiply the defining sum by \(t\), use the covariance relation, and
substitute \(b=ta\). Multiplication by \(t\) permutes \(F^\times\), so
\[
tM_j=\sum_a a^j e(ta)
    =\sum_b(t^{-1}b)^j e(b)=t^{-j}M_j.
\]
The change of variables is implemented by an explicit finite equivalence in
the Lean proof.
\end{proof}

\begin{corollary}[\lean{pow\_succ\_eq\_one\_of\_moment\_ne\_zero}]
\label{cor:moment-obstruction}
Under the same covariance relation,
\[
M_j(e)\neq0\quad\Longrightarrow\quad t^{j+1}=1.
\]
Without the nonvanishing hypothesis one has the factor identity
\[
\bigl(t^{j+1}-1\bigr)M_j(e)=0
\qquad(\lean{moment\_obstruction}).
\]
\end{corollary}

This is the kernel-checked Fourier argument behind the moment obstruction
\(\operatorname{im}_r(S)\subseteq T\) in the wider local theory. The present
file formalizes its algebraic engine; connecting the abstract function \(e\)
to the complete definition of \(S(p,r)\) remains outside this core.

\section{The golden moment}\label{sec:golden-moment}

Let \(\zeta\) satisfy \(\Phi_5(\zeta)=0\), and put
\[
 s=1+2(\zeta+\zeta^4),\qquad
 \eps=1+\zeta+\zeta^4,\qquad
 U_2=\prod_{a=1}^{4}(\zeta^a-1)^{a^2\bmod 5}.
\]
Thus \(s^2=5\), \(\eps=(1+s)/2\) whenever \(2\) is invertible, and
\[
U_2=(\zeta-1)(\zeta^2-1)^4(\zeta^3-1)^4(\zeta^4-1).
\]
The quintic index of \(U_2\) is the quadratic moment \(M_2\) in the
moment decomposition of the local Agrawal congruence.

\begin{theorem}[\lean{golden\_moment\_factorization}]
\label{thm:golden-factorization}
In every commutative ring in which
\(\zeta^4+\zeta^3+\zeta^2+\zeta+1=0\), one has
\[
\boxed{\quad U_2=s^5\eps^3.\quad}
\]
\end{theorem}

\begin{proof}[Certified proof sketch]
Set
\[
A=(\zeta-1)(\zeta^4-1),\qquad
B=(\zeta^2-1)(\zeta^3-1).
\]
Direct reductions modulo \(\Phi_5\), formalized without division, give
\[
AB=5,\qquad B=s\eps,\qquad s^2=5.
\]
Since \(U_2=AB^4=(AB)B^3\), it follows that
\(U_2=5(s\eps)^3=s^5\eps^3\).
\end{proof}

\begin{theorem}[Golden index theorem]
\label{thm:golden-index}
Let \(p\) be prime with \(5\mid p-1\), let \(\zeta\in\F_p\) satisfy
\(\Phi_5(\zeta)=0\), and suppose that the displayed elements are units.
For the discrete quintic index
\(\operatorname{ind}_5:\F_p^\times\to\Z/5\Z\) constructed in
\lean{quinticIndex},
\[
\boxed{\quad M_2=3\,\operatorname{ind}_5(\eps).\quad}
\]
The same identity holds for every multiplicative quintic character.
\par\smallskip\noindent
\textit{Lean declarations:}\quad
\lean{golden\_moment\_index}, \lean{zmod\_golden\_moment\_index}.
\end{theorem}

\begin{proof}[Certified proof sketch]
Every quintic index kills the fifth power \(s^5\), and sends
\(\eps^3\) to \(3\,\operatorname{ind}_5(\eps)\). Apply
Theorem~\ref{thm:golden-factorization}.
\end{proof}

\begin{remark}[Prior art and novelty boundary]\label{rem:golden-prior-art}
Williams and Hardy computed the quintic index of the golden unit in
Dickson coordinates \cite[Theorem~5]{WH}; that classical character is
not claimed here as new. The contribution of
Theorems~\ref{thm:golden-factorization}--\ref{thm:golden-index} is the
identification of the \emph{Agrawal quadratic moment} with that character.
We have not found this bridge in the targeted literature search, but
priority remains provisional pending specialist review.
\end{remark}

\section{The golden--cyclotomic entanglement}

\begin{theorem}[\lean{entanglement}]\label{thm:ent}
In any commutative ring, if
\(\zeta^4+\zeta^3+\zeta^2+\zeta+1=0\), then
\[
\zeta^3\,(1+\zeta+\zeta^4)^2\,(\zeta-1)^4 \;=\; 5 .
\]
In particular this holds in \(\Z[\zeta_5]\) with
\(\eps=1+\zeta+\zeta^4\) the image of the golden unit.
\end{theorem}

The proof is a single \lean{linear\_combination} with an explicit
cofactor; \(\zeta^5=1\) follows from the same hypothesis
(\lean{zeta\_pow\_five}).

\section{The Fermat-shadow glue}

\begin{lemma}[\lean{mul\_dvd\_gcd\_mul}]\label{lem:index}
If \(x\mid M\) and \(y\mid M\), then \(xy\mid\gcd(x,y)\,M\).
\end{lemma}

\begin{theorem}[\lean{fermat\_shadow}]\label{thm:shadow}
Let \(n\) be squarefree with \(5\nmid n\), and suppose that for every
prime \(q\mid n\)
\[
5^{\,n/q-1}\equiv1\pmod q .
\tag{\(\ast\)}
\]
Then \(n\mid 5^{\,n-1}-1\). Thus \(n\) satisfies the base-\(5\)
Fermat congruence; if \(n\) is composite, it is a Fermat pseudoprime
to base \(5\).
\end{theorem}

\begin{remark}
Hypothesis \((\ast)\) is the local norm constraint of the wider study;
in Theorem~\ref{thm:shadow} it is an explicit hypothesis. The next
theorem removes it: in our squarefree \(r=5\) setting, the Agrawal
congruence itself forces the base-\(5\) Fermat congruence.
\end{remark}

\begin{theorem}[\lean{agrawal\_fermat\_shadow}]\label{thm:bridge}
Let \(n\) be squarefree with \(5\nmid n\), and suppose the Agrawal
congruence at \(r=5\) holds:
\((X-1)^n\equiv X^n-1\pmod{n,\,X^5-1}\).
Then \(n\mid 5^{\,n-1}-1\). In particular, every squarefree composite
counterexample to Agrawal's conjecture at \(r=5\) prime to \(10\) is
a Fermat pseudoprime to base \(5\).
\end{theorem}

\begin{proof}[Certified proof sketch]
Reduce the congruence mod a prime \(q\mid n\); it becomes a
polynomial identity in \((\Z/q)[X]\), which can be evaluated at the
four elements \(z^j\) (\(j=1,\dots,4\)) of the quotient ring
\((\Z/q)[X]/\Phi_5\) (not a field in general), where \(z^5=1\) kills
the modulus. Multiplying the four evaluated identities and using the
explicit-cofactor identity \((z-1)(z^2-1)(z^3-1)(z^4-1)=5\)
(\lean{prod\_pow\_sub\_one}) gives \(5^{\,n}=5\) in the quotient
ring; injectivity of the constant embedding descends this equality
to \(\Z/q\), hence \(5^{\,n-1}=1\) there (\lean{agrawal\_local}),
and the squarefree gluing finishes. No Frobenius descent and no norm
maps are needed.
\end{proof}

\begin{remark}[Prior art]
Lenstra observed the more general implication
\(r^n\equiv r\pmod n\) directly from Agrawal's congruence
\cite[Remark~5]{LPproblems}. Theorem~\ref{thm:bridge} is therefore not
claimed as new mathematics: its contribution here is an end-to-end Lean
formalization of the squarefree \(r=5\) specialization.
\end{remark}

\section{The two-adic jaw}

The bridge constrains a squarefree solution of the Agrawal congruence
only through the multiplicative order of \(5\). On the primes that are
inert in \(\mathbf{Q}(\sqrt5)\) that constraint is rigid at the
prime \(2\).

\begin{theorem}[\lean{agrawal\_two\_adic\_jaw}]\label{thm:jaw}
Let \(n\) be squarefree with \(5\nmid n\) and suppose the Agrawal
congruence at \(r=5\) holds. Let \(q\mid n\) be an odd prime with
\(q\equiv\pm2\pmod5\). Then
\[
v_2(q-1)\;\le\;v_2(n-1).
\]
\end{theorem}

\begin{proof}[Certified proof sketch]
By Theorem~\ref{thm:bridge}, \(\operatorname{ord}_q(5)\) divides
\(n-1\). Since \(q\equiv\pm2\pmod5\), quadratic reciprocity gives
\((5/q)=-1\) (\lean{not\_isSquare\_five}), so \(5^{(q-1)/2}=-1\)
(\lean{euler\_nonsquare}) and the order does \emph{not} divide
\((q-1)/2\). A divisor of \(2^{a}u\) with \(u\) odd that fails to
divide \(2^{a-1}u\) is 2-adically saturated
(\lean{pow\_two\_dvd\_of\_not\_dvd\_half}), whence
\(v_2(\operatorname{ord}_q5)=v_2(q-1)\), and the claim follows; the
local form is \lean{two\_adic\_jaw\_local}.
\end{proof}

\begin{corollary}[\lean{agrawal\_inert\_mod\_four}]\label{cor:modfour}
Under the same hypotheses, if \(n\equiv3\pmod4\) then every prime
factor of \(n\) that is inert in \(\mathbf{Q}(\sqrt5)\) satisfies
\(q\equiv3\pmod4\).
\end{corollary}

\begin{remark}
The split primes behave in the opposite way: for \(q\equiv\pm1\pmod5\)
one has \((5/q)=+1\), hence \(\operatorname{ord}_q(5)\mid(q-1)/2\) and
\(v_2(\operatorname{ord}_q5)<v_2(q-1)\), so no constraint is
transmitted. Both statements above are unconditional.
\end{remark}

\section{The bounded order}

The bridge and the jaw both read the congruence through the
multiplicative order of an element. The element that carries the
congruence itself is \(\zeta-1\), where \(\zeta\) is the class of
\(X\) in \(T=\mathbf{F}_p[X]/\Phi_5\). A priori its order is only
known to divide \(|T^\times|=p^4-1\). It divides much less.

\begin{theorem}[\lean{order\_bounded}]\label{thm:order}
Let \(p\) be a prime with \(p\equiv\pm2\pmod5\), so that \(\Phi_5\)
stays irreducible modulo \(p\), and let \(\zeta\) be the class of
\(X\) in \(T=\mathbf{F}_p[X]/\Phi_5\). Then
\[
(\zeta-1)^{10p^2}=(\zeta-1)^{10}.
\]
\end{theorem}

The displayed identity is the statement checked by the cited Lean
declaration. Mathematically, irreducibility makes \(T\) a field and
\(\zeta-1\neq0\); cancellation therefore gives the finite-field
corollary that the multiplicative order of \(\zeta-1\) divides
\(10(p^2-1)\). This order formulation is not bundled into
\lean{order\_bounded}.

\begin{proof}[Certified proof sketch]
Write \(u=\zeta-1\). In characteristic \(p\) the Frobenius is
additive, so \(u^p=\zeta^p-1\), and since \(\zeta^5=1\) the exponent
only matters modulo \(5\) (\lean{zeta5\_pow\_mod}), giving
\(u^p=\zeta^{p\bmod 5}-1\) (\lean{zeta\_sub\_one\_pow\_p}). Applying
the Frobenius once more multiplies that exponent by \(p\) again, and
here the two inert classes merge: \(2^2\equiv3^2\equiv4\pmod5\), so in
both cases \(u^{p^2}=\zeta^4-1\)
(\lean{zeta\_sub\_one\_pow\_sq}). Finally \(\zeta^4=\zeta^{-1}\)
turns that into
\[
u^{p^2}=\zeta^{-1}-1=-(\zeta-1)\zeta^{-1}=-\zeta^4\,u
\]
(\lean{order\_bound\_relation}): the square of the Frobenius merely
scales \(u\) by \(-\zeta^4\). Raising to the tenth power kills the
scalar, since \((-1)^{10}=1\) and \(\zeta^{40}=(\zeta^5)^8=1\).
\end{proof}

\begin{remark}
The statement is phrased as \(u^{10p^2}=u^{10}\) rather than as a
divisibility of the order, so that no invertibility of \(u\) is
needed. It is the form actually used downstream: the factor \(10\) is
what forces the modulus \(80\) to appear in the congruence conditions
on the prime factors of a putative counterexample, and it is what
makes an exact-order computation factor a ten-digit integer instead of
a twenty-digit one.
\end{remark}

\section{An elementary box}

The following identity is used to confine the largest prime factor of
a three-factor solution to an interval. It is stated over any
commutative ring and proved by expansion.

\begin{theorem}[Box identity]\label{thm:box}
In any commutative ring, for all \(p,q,r\),
\[
\begin{split}
&2\bigl[(pq-1)(pr-1)(qr-1)
        -(p^2-1)(q^2-1)(r^2-1)\bigr]\\
&\qquad=(p^2-1)(q-r)^2+(q^2-1)(p-r)^2
        +(r^2-1)(p-q)^2 .
\end{split}
\]
In particular the left-hand side is non-negative whenever
\(p,q,r\ge1\).
\par\smallskip\noindent
\textit{Lean declarations:}\par\noindent
\lean{prod\_pairs\_sub\_prod\_squares},\par\noindent
\lean{prod\_pairs\_ge\_prod\_squares}.
\end{theorem}

\noindent Combined with \lean{lt\_two\_mul\_of\_sq\_le}, which turns
\(r^2\le 2I(pq)\) with \(q<r\) into \(r<2Ip\), this bounds the third
prime by an explicit multiple of the smallest one.

\section{The Lenstra--Pomerance proposition}

The statement below is the one on which the search for counterexamples at
\(r=5\) rests. It appears in the AIM notes of Lenstra and Pomerance
\cite[pp.~30--32]{LPremarks} and is quoted throughout the subsequent
literature; to our knowledge it had not been machine-checked before.

Mathlib does not currently contain Carmichael numbers or Korselt's criterion,
so those had to be introduced: \lean{IsCarmichael} and
\lean{IsLucasCarmichael} package the two divisibility conditions used below.
We also prove that a Carmichael number is a Fermat pseudoprime to every
coprime base (declaration
\texttt{IsCarmichael.}\allowbreak\texttt{fermatPsp}). This implication is
\emph{not} used in Theorem~\ref{thm:lenstra}, which needs only the
divisibility conditions themselves.

\begin{theorem}[\lean{lenstra\_proposition\_card}]\label{thm:lenstra}
Let \(n\) be squarefree, suppose the number \(k\) of its prime factors
satisfies \(k\equiv1\pmod4\), suppose every prime \(p\mid n\) satisfies
\(p\equiv3\pmod{80}\), and suppose Korselt's two conditions hold:
\((p-1)\mid(n-1)\) and \((p+1)\mid(n+1)\) for every \(p\mid n\). Then
\[
(X-1)^n\equiv X^n-1 \pmod{n,\ X^5-1}
\qquad\text{and}\qquad
n^2\not\equiv1\pmod5 .
\]
\end{theorem}

These are the hypotheses of the original statement. That \(n\equiv3\pmod{80}\)
follows from them is itself proved (\lean{mod\_eighty\_of\_card}): for
squarefree \(n\) the product of the prime factors is \(n\), each factor is
\(3\) modulo 80, and \(3^4=81\equiv1\pmod{80}\) reduces the exponent to the
class of \(k\) modulo 4. A variant assuming \(n\equiv3\pmod{80}\) directly
is available as \lean{lenstra\_proposition}.

The second conclusion is what makes such an \(n\), \emph{if composite}, a
genuine counterexample: Agrawal's conjecture allows the escape
\(n^2\equiv1\pmod r\), and here it is closed
(\lean{sq\_not\_one\_mod\_five}). The proviso matters: primes satisfy the
hypotheses trivially---\(n=83\) is one---and are of course not
counterexamples. Only a composite witness would be one, and none is known.

\par\smallskip
\begin{proof}[Certified proof sketch]
The route is the one of the original proof, which already contains the
identity \((\zeta_5-1)^{p^2}=-\zeta_5^{-1}(\zeta_5-1)\), the resulting bound
on the order, and the reduction to \(n\equiv p\); we repackage it through a
single observation: \emph{the three conditions imposed on
\(n\) are satisfied by \(p\) itself}, since \(p\equiv1\pmod{p-1}\),
\(p\equiv-1\pmod{p+1}\) and \(p\equiv3\pmod{80}\). Hence
\[
n\equiv p \pmod{\operatorname{lcm}(p-1,\,p+1,\,80)},
\qquad
\operatorname{lcm}(p-1,\,p+1,\,80)=10(p^2-1)
\]
(\lean{lcm\_three\_eq}, \lean{sub\_dvd\_of\_korselt}), and that modulus is
exactly the one bounding the order of \(\zeta-1\) by
Theorem~\ref{thm:order}. Therefore \((\zeta-1)^n=(\zeta-1)^p\), and the
Frobenius gives \((\zeta-1)^p=\zeta^p-1=\zeta^3-1=\zeta^n-1\), which is the
congruence in the cyclotomic component (\lean{lenstra\_local}). On the
component \(X\mapsto1\) both sides vanish, and since
\(X^5-1=(X-1)\Phi_5\) with the two factors coprime for \(p\neq5\), the two
pieces glue to the congruence modulo \(p\) (\lean{agrawal\_mod\_p}, using
\lean{dvd\_of\_dvd\_both}). Finally, for squarefree \(n\) a congruence
holding modulo every prime factor holds modulo \(n\): dividing by the monic
\(X^5-1\) over \(\mathbf Z\), every coefficient of the remainder is divisible
by each \(p\mid n\), hence by \(n\) (\lean{congruence\_of\_local}).
\end{proof}

\begin{remark}
Nothing here asserts that a composite such \(n\) exists. A composite number
satisfying the hypotheses would be simultaneously a Carmichael and a
Lucas--Carmichael number; no such integer is known. Williams posed this
existence problem in 1977 \cite{Williams1977,McIntoshDipra2015}; Pomerance's 1984
Baillie--PSW heuristic later imposed the same paired Korselt conditions
\cite{Pomerance1984}. Note also that the proposition gives a \emph{sufficient}
condition for producing a counterexample, not a necessary one: it does not say
where all counterexamples must lie. It certifies the tool used to look for
them; it does not decide Agrawal's conjecture in either direction.
\end{remark}

\section{The forced partition}

The simultaneous Carmichael and Lucas--Carmichael conditions impose a
prime-by-prime separation on the factors of any hypothetical witness.

\begin{theorem}[\lean{common\_divisor\_forced}]\label{thm:partition}
Let \(n=pm=p'm'\), and suppose that both Korselt divisibilities hold
at \(p\) and at \(p'\):
\[
p-1\mid n-1,\quad p+1\mid n+1,\qquad
p'-1\mid n-1,\quad p'+1\mid n+1.
\]
If \(d\mid p-1\) and \(d\mid p'+1\), then \(d\mid2\). Consequently, for an
odd prime \(q\), the two sets
\[
\{p\mid n:q\mid p-1\},\qquad
\{p\mid n:q\mid p+1\}
\]
cannot both be nonempty (\lean{partition\_forced}).
\end{theorem}

\begin{proof}[Certified proof sketch]
Writing \(n=pm\), the two divisibilities at \(p\) give
\(p-1\mid m-1\) and \(p+1\mid m-1\)
(\lean{dvd\_sub\_of\_korselt\_pair}). Since \(d\mid p-1\), it follows
that \(d\mid m-1\), and therefore \(n\equiv p\pmod d\). The same
argument at \(p'\), using \(d\mid p'+1\), gives
\(n\equiv p'\pmod d\), so \(p\equiv p'\pmod d\). But the two side
conditions give \(p\equiv1\pmod d\) and \(p'\equiv-1\pmod d\).
Thus \(d\mid2\).
\end{proof}

\begin{remark}[Scope and prior art]
The compatibility condition
\(\gcd(p_i-1,p_j+1)=2\) for Williams sets is stated explicitly by
McIntosh \cite[\S4]{McIntosh2014}; it is not a new mathematical claim of
this paper. The contribution here is a kernel-checked arbitrary-divisor
form and its reusable corollaries. The intermediate congruence
\(n\equiv p\pmod{(p^2-1)/2}\), and the stronger consequence that all
prime factors of a simultaneous Carmichael and Lucas--Carmichael number
are congruent modulo \(12\), already appear in Leng's 2017 study
\cite{Leng2017}; the named mod-\(3\) corollary is
\lean{three\_congruence\_forced}. Theorem~\ref{thm:partition} is
a useful structural sieve, not a proof that simultaneous numbers do not
exist. In particular, the exclusions contributed by different \(q\)
need not be statistically independent.

There is also no contradiction from size alone. If
\(L=\operatorname{lcm}_{p\mid n}(p-1)\) and
\(M=\operatorname{lcm}_{p\mid n}(p+1)\), then \(LM\mid n^2-1\), but
\[
 LM\le\prod_{p\mid n}(p^2-1)<n^2
\]
is automatic. The non-tautological information is instead at each prime-power
level: if \(I^-_{\ell,e}=\{i:\ell^e\mid p_i-1\}\), then
\(\prod_{i\notin I^-_{\ell,e}}p_i\equiv1\pmod{\ell^e}\); if
\(I^+_{\ell,e}=\{i:\ell^e\mid p_i+1\}\), then
\(\prod_{i\notin I^+_{\ell,e}}p_i
\equiv(-1)^{|I^+_{\ell,e}|+1}\pmod{\ell^e}\).
These incidence congruences, rather than the global inequality, are the
next structural sieve.

The cubic parametrization used in the five-factor search does not provide
an independent constraint: the strong complementary row
\(10(p^2-1)\mid m-1\) is equivalent, for \(p\neq0\), to
\[
 pm=p+10u(p^3-p)
\]
for some \(u\). Both directions are kernel-checked:
\begin{center}
\small
\lean{cubic\_signature\_of\_strong\_row}\\
\lean{strong\_row\_of\_cubic\_signature}.
\end{center}
\end{remark}

\section{Computational certificates}

Independently of the kernel-checked statements above, the repository
carries replayable certificates, each embedding its own detection
criterion and reconstruction data, with SHA-256 manifests:

\begin{enumerate}
\item seven \emph{certified-empty fibres} of a three-factor sieve
associated with the study (pairs
\((13,37)\), \((17,113)\), \((13,277)\), \((23,167)\), \((7,823)\),
\((83,347)\), \((463,1327)\)): for each pair, the finitely many
levels required by the certificate are integrally factored, all prime
factors are proven prime, and no prime in the certificate's detector
classes survives;
\item a self-certifying census of \(H_n\) for all
\(n\equiv4\pmod5\), \(n\le10^5\): \(9{,}725\) nontrivial values,
all factorizations proven, and \emph{no} prime factor
\(\equiv\pm1\pmod5\) occurs.
\end{enumerate}

These are certificates of finite computations; they prove exactly
what they state and nothing beyond, and both replay from the
repository alone. Separately, the repository freezes as
\emph{hash-only provenance} (hashes, not artifacts) a SHA-256 index
of the working material of the study, including a prime-by-prime
corpus up to \(10^9\); nothing in this note relies on it.

\section{Two open problems}

We pose the two statements that our working material identifies as
the central open problems of the \(r=5\) study. We claim no proof and
no partial proof here.

\begin{conjecture}[golden inertia]\label{conj:inertia}
For every \(n\equiv4\pmod5\), every odd prime factor of \(H_n\)
distinct from \(5\) is \(\equiv\pm2\pmod5\).
\end{conjecture}

Theorem~\ref{thm:support} shows that every prime
\(p\equiv2\pmod5\) does occur in some \(H_{n}\) with
\(n\equiv4\pmod5\); the census in item (2) above verifies
Conjecture~\ref{conj:inertia} for all \(n\le10^5\).

\begin{conjecture}[fibre emptiness]\label{conj:fibres}
The pattern of item (1) above is general: for every admissible pair,
the corresponding fibre is empty.
\end{conjecture}

\section{Formalization data}

At the release state audited for this version there are twenty-five modules
and \(2682\) source lines; no \lean{sorry}, no \lean{admit}, and no axioms
beyond Mathlib's. Mathlib is
pinned to a fixed upstream commit; a clean clone builds with
\lean{lake exe cache get \&\& lake build --wfail}. The table mapping each
statement of this note to its declaration and file is in the
repository \lean{README}.

\section{Reproducibility map}

The repository is organized so that the three kinds of evidence never
collapse into one another:
\begin{enumerate}
\item the kernel layer is rebuilt with
  \lean{lake build --wfail}; its flagship modules cover the moment
  obstruction, the golden moment, and the forced partition;
\item the finite-computation layer is replayed by running
  \path{tools/verify_certificates.py}; the optional \lean{--full} flag
  recomputes the complete \(H_n\) census;
\item the literature layer is recorded claim by claim in
  \path{NOVELTY_AND_PRIOR_ART.md}, including negative novelty findings.
\end{enumerate}
A read-only GitHub Actions workflow runs the first two checks with Lean
warnings promoted to errors. The release gate remains private until a human
specialist has reviewed the golden-moment priority claim.

\begin{thebibliography}{9}

\bibitem{AKS}
M.~Agrawal, N.~Kayal, and N.~Saxena,
\emph{PRIMES is in P},
Ann. of Math. 160 (2004), 781--793,
\url{https://doi.org/10.4007/annals.2004.160.781}.

\bibitem{LPremarks}
H.~W. Lenstra, Jr. and C.~Pomerance,
\emph{Remarks on Agrawal's conjecture},
AIM note, 2003,
\url{https://aimath.org/WWN/primesinp/articles/html/50a/}.

\bibitem{LPproblems}
H.~W. Lenstra, Jr. and C.~Pomerance,
\emph{Problems concerning Agrawal's conjecture},
AIM note, 2003,
\url{https://aimath.org/WWN/primesinp/articles/html/38a/}.

\bibitem{LPgauss}
H.~W. Lenstra, Jr. and C.~Pomerance,
\emph{Primality testing with Gaussian periods},
J. Eur. Math. Soc. 21 (2019), 1229--1269,
\url{https://doi.org/10.4171/JEMS/861}.

\bibitem{Vana}
T.~V\'a\v na,
\emph{Agrawal's conjecture and Carmichael numbers},
Student Scientific Conference, Comenius University Bratislava, 2009,
\url{https://jcmf.cz/sites/default/files/cms/archiv_svoc/2009/prace/7/Vana.pdf}.

\bibitem{WH}
K.~S. Williams and K.~Hardy,
\emph{A congruence for the index of a unit of a real abelian number field},
Acta Arith. 46 (1985), 57--72,
\url{https://eudml.org/doc/205987}.

\bibitem{Pomerance1984}
C.~Pomerance,
\emph{Are there counter-examples to the Baillie--PSW primality test?},
1984,
\url{https://math.dartmouth.edu/~carlp/dopo.pdf}.

\bibitem{Williams1977}
H.~C. Williams,
\emph{On numbers analogous to the Carmichael numbers},
Canad. Math. Bull. 20 (1977), 133--143,
\url{https://doi.org/10.4153/CMB-1977-019-8}.

\bibitem{McIntosh2014}
R.~J. McIntosh,
\emph{Carmichael numbers with \(p+1\mid n-1\)},
Integers 14 (2014), Paper A59,
\url{https://ftp.icm.edu.pl/packages/EMIS/journals/INTEGERS/papers/o59/o59.pdf}.

\bibitem{McIntoshDipra2015}
R.~J. McIntosh and S.~Mitra Dipra,
\emph{Carmichael numbers with \(p+1\mid n+1\)},
J. Number Theory 147 (2015), 81--91,
\url{https://doi.org/10.1016/j.jnt.2014.06.005}.

\bibitem{Leng2017}
A.~Leng,
\emph{Independence of the Miller--Rabin and Lucas probable prime tests},
2017,
\url{https://math.mit.edu/research/highschool/primes/materials/2016/Leng.pdf}.

\end{thebibliography}

\end{document}
